\section{Conclusions} \label{sec:conclusions}

In this work, we have presented a CRL framework for whole-body loco-manipulation for quadrupedal robots in lunar surface operations. By unifying locomotion and manipulation into a single end-to-end policy, our approach overcomes the limitations of decoupled controllers, allowing the robot to effectively coordinate base mobility with arm dexterity to track 6D task-space commands.

Our simulation campaign shows that the proposed framework achieves mean position errors of $3.22\,$cm and mean orientation errors of $6.83^{\circ}$ under nominal lunar-analogue conditions, while combining locomotion and manipulation in a single whole-body control policy. This precision is retained across terrain morphologies and regolith properties over the operationally representative difficulty range, degrading to approximately $13\,$cm only on the most extreme crater slopes evaluated ($30^{\circ}$), which lie at the limit of what a platform of this class can traverse at all. These errors are of the same order as those reported by state-of-the-art terrestrial whole-body tracking methods, obtained here with a single controller and under lunar gravity; we note in \cref{exp:tracking} that this is not a controlled comparison, as platform, task distribution and environment all differ. Crucially, the integration of the CaT mechanism allowed us to embed mission-critical safety requirements, such as stability margins, impact force limits, and power efficiency, directly into the learning process. The experimental results show that this formulation reduces soft constraint violation time by up to an order of magnitude relative to an unconstrained baseline, with no hard constraint violations observed over the evaluated episodes.

The policy also absorbed domain randomization and to contact dynamics not encountered during training. It maintained stability and low tracking error across varying terrain morphologies and diverse regolith properties (loose to dense soil), which supports, without establishing, its suitability for the conditions expected on the lunar surface. By addressing the specific challenges of reduced gravity and compliant regolith contact, this work provides a quantitative simulation baseline against which subsequent hardware studies can be measured.

Several directions follow from this work. The most immediate is validation on physical hardware under gravity offload, together with a contact model capable of representing regolith behaviour under sustained loading, which would allow the constraint set to be evaluated against physically grounded failure modes. Controlled comparisons against decoupled and unconstrained architectures would quantify the benefit of unified whole-body control that our task formulation motivates but does not measure. Beyond these, coupling the learned policy to a certified safety filter offers a route to the assurance that probabilistic constraint enforcement alone cannot provide, and extending the framework to contact-rich operations would bring it closer to the sampling, assembly and maintenance tasks that lunar surface missions will require.
